%! Modeling with Exponential Functions: Interest, Half-Life, and e — Unit 6, Lesson 6.4 — Group Activity (Answer Key)
\documentclass[10pt]{article}
\usepackage{algebra2-article}
\usepackage{algebra2-key}   % pulls in -boxes; do NOT also load -boxes
\newcommand{\TallMath}[1]{$\displaystyle #1\rule[-1.4em]{0pt}{3.2em}$}

\begin{document}

\pageheader{Unit 6, Lesson 6.4}{Group Activity --- Answer Key}
\namepartnerperiod

\vspace{0.08in}

\begin{headlinebox}{forestbg}
\small\textbf{The Base Is One Period. The Exponent Counts Periods.} Every model below is
$y=ab^{x}$ wearing a costume. Before either partner touches a calculator, both of you must agree out
loud on two things: \textbf{what is one period here?} and \textbf{how many of them fit in the time
given?} Answer those and the substitution writes itself. Round money to the nearest cent and
milligrams to two decimal places.
\end{headlinebox}

\vspace{0.06in}

\begin{tcolorbox}[colback=white, colframe=black!40, title=\textbf{Tier R --- Remediate},
    fonttitle=\bfseries, arc=1.5mm, left=3mm, right=3mm, top=2mm, bottom=2mm, breakable]
\textbf{Part 1 --- Set up before you solve.} \$2000 is invested at $3\%$ for $10$ years. Fill in one
row at a time; the last column is the whole point.
\begin{center}
\small
\renewcommand{\arraystretch}{1.6}
\begin{tabular}{|l|c|c|c|c|}
\hline
\textbf{Compounded} & $n$ & \textbf{base} $1+\frac{r}{n}$ & \textbf{exponent} $nt$ & \textbf{$A=2000\left(1+\frac{r}{n}\right)^{nt}$} \\
\hline
annually  & $1$ & $1.03$ & $10$ & \ans{\$2687.83} \\
\hline
quarterly & $4$ & \ans{$1.0075$} & \ans{$40$} & \ans{\$2696.70} \\
\hline
\end{tabular}
\end{center}

\vspace{2pt}
Compounding four times as often earned an extra \ans{\$8.87} over ten years. Was that more or less
than your partner expected? \ans{Answers vary --- almost everyone guesses high.}

\vspace{6pt}
\textbf{Part 2 --- A medicine that halves.} A $600$ mg dose has a half-life of $12$ hours.

\vspace{3pt}
\hspace*{1.0em} Model: $A(t)=$ \ans{$600\left(\frac12\right)^{t/12}$}, \; where $t$ is measured in \ans{hours}.

\vspace{3pt}
\begin{center}
\small
\renewcommand{\arraystretch}{1.4}
\begin{tabular}{|l|c|c|c|c|c|}
\hline
$t$ (hours) & $0$ & $12$ & $24$ & $36$ & $48$ \\ \hline
halvings $t/12$ & $0$ & \ans{1} & \ans{2} & \ans{3} & \ans{4} \\ \hline
$A(t)$ (mg) & $600$ & \ans{300} & \ans{150} & \ans{75} & \ans{37.5} \\ \hline
\end{tabular}
\end{center}

\vspace{2pt}
Every row cuts the one before it in half, yet the amount never reaches $0$. Which feature of the graph
says so? \ans{the horizontal asymptote at $A=0$ --- the range is $(0,600]$}

\vspace{5pt}
\textbf{Part 3 --- One period, named.} For each model, say what one period is and how many periods fit
in the time asked for.

\vspace{2pt}
\begin{center}
\small
\renewcommand{\arraystretch}{1.6}
\begin{tabularx}{\linewidth}{@{}l X X@{}}
\hline
\textbf{Model} & \textbf{One period is\dots} & \textbf{Periods in the time given} \\
\hline
$4000(1.02)^{4t}$, \; $t=3$ yr & \ans{one quarter} & \ans{$12$} \\
$50\left(\frac12\right)^{t/9}$, \; $t=27$ days & \ans{$9$ days (one half-life)} & \ans{$3$} \\
$300\cdot2^{\,t/4}$, \; $t=20$ hours & \ans{$4$ hours (one doubling)} & \ans{$5$} \\
\hline
\end{tabularx}
\end{center}
\end{tcolorbox}

\begin{tcolorbox}[colback=white, colframe=black!40, title=\textbf{Tier A --- Approaching Proficiency},
    fonttitle=\bfseries, arc=1.5mm, left=3mm, right=3mm, top=2mm, bottom=2mm, breakable]
\textbf{Part 1 --- Two offers, one deposit.} A credit union offers \$5000 at $3.6\%$ compounded
\emph{monthly}. A bank offers the same \$5000 at $3.6\%$ compounded \emph{continuously}. Both are
left alone for $8$ years.

\vspace{3pt}
\begin{center}
\small
\renewcommand{\arraystretch}{1.7}
\begin{tabularx}{\linewidth}{@{}l X X c@{}}
\hline
& \textbf{Formula used} & \textbf{Substitution} & \textbf{Balance} \\
\hline
Credit union & \ans{$P\left(1+\frac{r}{n}\right)^{nt}$} & \ans{$5000(1.003)^{96}$} & \ans{\$6665.91} \\
Bank & \ans{$Pe^{\,rt}$} & \ans{$5000e^{\,0.288}$} & \ans{\$6668.79} \\
\hline
\end{tabularx}
\end{center}

\vspace{2pt}
\textbf{(a)} Which is better, and by how much? \ans{the bank, by \$2.88 over eight years}

\vspace{3pt}
\textbf{(b)} A classmate says continuous compounding should ``obviously'' win by a lot, since it
compounds infinitely often. Using your two numbers, explain what is actually true.
\ansline{Continuous compounding is the ceiling, but it is a very low ceiling: on \$5000 over eight years it beats monthly by \$2.88, about $0.04\%$.}
\ansline{Each extra chop of the year adds less than the one before, so ``infinitely often'' is worth barely more than ``twelve times a year.''}

\vspace{5pt}
\textbf{Part 2 --- Between the halvings.} A $120$ mg dose has a half-life of $5$ hours.

\vspace{2pt}
\hspace*{1.0em} $A(t)=$ \ans{$120\left(\frac12\right)^{t/5}$} \hfill $A(10)=$ \ans{30 mg} \hfill $A(17)=$ \ans{11.37 mg}

\vspace{3pt}
$A(17)$ is not a whole number of halvings. Explain why that is fine, and check that your answer is the
size it ought to be by naming the two table values it must fall between.
\ansline{$17\div5=3.4$ halvings --- a fractional exponent is a legal exponent (Lesson 5.1), and $t$ does not have to land on a half-life.}
\ansline{It must sit between $A(15)=15$ mg and $A(20)=7.5$ mg, and $11.37$ does.}

\vspace{5pt}
\textbf{Part 3 --- Error analysis.} An account pays $8\%$ compounded quarterly. Three students model
\$1000 after $5$ years. Name each error, then write the correct model and value.

\vspace{2pt}
\textbf{(a)} Renata writes $A=1000(1.08)^{20}$.
\ansline{She counted the periods correctly ($20$ quarters) but never split the rate --- her base $1.08$ pays a full year's interest \emph{every quarter}.}
\ansline{Her answer, \$4660.96, is the balance after $20$ \emph{years}, not $5$. Base and exponent must agree on what one period is.}

\vspace{1pt}
\textbf{(b)} Miles writes $A=1000(1.02)^{5}$.
\ansline{He split the rate correctly ($1.02$ per quarter) but counted years instead of quarters, so his model only runs for five quarters.}
\ansline{His answer, \$1104.08, is the balance after $1\frac14$ years --- the opposite half of Renata's mistake.}

\vspace{1pt}
\textbf{(c)} Correct model: $A=$ \ans{$1000(1.02)^{20}$} $=$ \ans{\$1485.95}

\vspace{3pt}
Renata and Miles made \emph{opposite} halves of the same mistake. Say what the one sentence is that
both of them broke. \ans{The base is what happens in one period; the exponent counts periods.}
\end{tcolorbox}

\begin{tcolorbox}[colback=white, colframe=black!40, title=\textbf{Tier E --- Extension},
    fonttitle=\bfseries, arc=1.5mm, left=3mm, right=3mm, top=2mm, bottom=2mm, breakable]
\textbf{Part 1 --- Where the chopping stops.} You saw $\left(1+\frac1m\right)^{m}$ close in on $e$.
Now watch the same thing happen at a realistic rate: $\left(1+\frac{0.06}{n}\right)^{n}$, the growth
factor for one year at $6\%$.

\vspace{3pt}
\begin{center}
\small
\renewcommand{\arraystretch}{1.4}
\begin{tabular}{|l|c|c|c|c|c|}
\hline
$n$ & $1$ & $4$ & $12$ & $365$ & continuous \\
\hline
one-year factor & $1.060000$ & $1.061364$ & $1.061678$ & $1.061831$ & \ans{$1.061837$} \\
\hline
\end{tabular}
\end{center}
\hspace*{1.0em} The last cell is $e^{\,0.06}$. Compute it to six decimal places: \ans{$1.061837$}

\vspace{3pt}
\textbf{(a)} A bank advertises ``compounded every \emph{second}.'' How much more than daily
compounding could that possibly pay on \$1000 in one year? \ans{about half a cent (\$0.005)}

\vspace{3pt}
\textbf{(b)} Why is $e^{\,0.06}$ a \emph{ceiling} rather than just the next entry in the table?
\ansline{No value of $n$ ever reaches it --- the factors climb toward $e^{0.06}$ and stop getting closer in any way you could notice.}
\ansline{``Every second,'' ``every instant,'' and every faster schedule are all still below it, so nothing beyond continuous compounding exists to try.}

\vspace{6pt}
\textbf{Part 2 --- The Rule of 72.} Bankers estimate doubling time with a shortcut: at $r\%$ per year
compounded annually, money doubles in about $\frac{72}{r}$ years. Test it against the graph-reading you
did in the notes.

\vspace{3pt}
\begin{center}
\small
\renewcommand{\arraystretch}{1.6}
\begin{tabularx}{\linewidth}{@{}c X c X c@{}}
\hline
\textbf{Rate} & \textbf{Rule of 72 estimate} & \textbf{$(1+r)^{t}$ at $t=$ estimate} & \textbf{Over or under 2?} & \textbf{True doubling time} \\
\hline
$5\%$ & \ans{$14.4$ yr} & \ans{$2.019$} & \ans{over} & about $14.2$ yr \\
$8\%$ & \ans{$9$ yr} & \ans{$1.999$} & \ans{under} & about $9.0$ yr \\
$12\%$ & \ans{$6$ yr} & \ans{$1.974$} & \ans{under} & about $6.1$ yr \\
\hline
\end{tabularx}
\end{center}

\vspace{2pt}
\textbf{(a)} At which of the three rates is the shortcut most accurate? \ans{at $8\%$ --- it is off by about a week}

\vspace{2pt}
\textbf{(b)} A shortcut is only worth having if you know when to distrust it. Based on your table, when
does the Rule of 72 start to go wrong?
\ansline{It drifts as you move away from $8\%$: at $5\%$ it overshoots the doubling and at $12\%$ it falls short, and the error grows at the high end.}
\ansline{It is an estimate for ordinary interest rates; the further from about $8\%$ you go, the less you should trust the last digit.}

\vspace{6pt}
\textbf{Part 3 --- Which of these can you actually finish?} Each question below reduces to an
exponential equation. Reduce it, then decide whether Lesson 6.3's common-base method finishes the job.

\vspace{3pt}
\begin{center}
\small
\renewcommand{\arraystretch}{1.8}
\begin{tabularx}{\linewidth}{@{}l X c X@{}}
\hline
\textbf{Question} & \textbf{Equation, after dividing} & \textbf{Common base?} & \textbf{If yes, the answer} \\
\hline
$600$ mg, $h=12$ h: when is $75$ mg left? & \ans{$\left(\frac12\right)^{t/12}=\frac18$} & \ans{yes} & \ans{$t=36$ h} \\
$300$ cells, $d=4$ h: when are there $2400$? & \ans{$2^{\,t/4}=8$} & \ans{yes} & \ans{$t=12$ h} \\
\$2000 at $3\%$ annually: when is it \$4000? & \ans{$(1.03)^{t}=2$} & \ans{no} & \ans{---} \\
$600$ mg, $h=12$ h: when is $200$ mg left? & \ans{$\left(\frac12\right)^{t/12}=\frac13$} & \ans{no} & \ans{---} \\
\hline
\end{tabularx}
\end{center}

\vspace{2pt}
\textbf{(a)} Two of the four defeated the method. For each one, say whether the answer \emph{exists},
and how you know without solving.
\ansline{Both answers exist. $2$ and $\frac13$ are positive, so each sits in the range $(0,\infty)$ and the horizontal line really does cross the curve.}
\ansline{Only the \emph{method} failed: neither $2$ nor $\frac13$ is a power of the base. Compare $2^{x}=-3$, which has no solution at all.}

\vspace{2pt}
\textbf{(b)} Rows 1 and 4 are the same drug with the same half-life. One works and one does not. What
is it about the \emph{target number} that decides it?
\ansline{$75$ is $600$ halved three times, so the target is a power of $\frac12$ times the start; $200$ is a \emph{third} of $600$, and no number of halvings produces a third.}
\ansline{Common-base solving only reaches targets the base can actually land on: $\frac12$, $\frac14$, $\frac18$, $\frac1{16}$ of the start --- nothing in between.}
\end{tcolorbox}

\begin{teachernote}
\textbf{Tier R Part 3 is the diagnostic, and it takes ninety seconds to read.} A student who can name
``one period'' for all three rows will get every computation on this page right; a student who cannot
will get all of them wrong in the same way. Check Part 3 before you check Part 1 --- it tells you
whether to re-teach or to let them work.

\textbf{Tier A Part 3 is written so the two errors cannot be confused.} Renata keeps the annual rate
with the quarterly count; Miles keeps the quarterly rate with the annual count. They need different
corrections, and the closing question is the one that matters --- both students broke the same sentence
in opposite directions. If a group writes ``they both made a mistake with the exponent,'' push back:
only one of them did.

\textbf{Tier A Part 1(b) is worth a whole-class share-out.} The $\$2.88$ result over eight years is
genuinely counterintuitive and it is the honest answer to the Hook. Continuous compounding matters
because it is the \emph{model} for things that change without pausing, not because it makes anyone
richer.

\textbf{Tier E Part 2 is not a formula to memorize.} The point is calibration: the Rule of 72 is
excellent near $8\%$ and drifts either side of it. Students who ask \emph{why} $72$ are asking a
logarithm question --- tell them so, and tell them it is Unit 7. Do not derive it.

\textbf{Tier E Part 3 is the closing beat of the lesson and should be debriefed even if groups do not
finish it.} Rows 1 and 4 are the whole idea: the same model, the same half-life, and one target that
common-base solving reaches and one it cannot. Part (a) is where students must say \emph{the answer
exists, the method failed} --- the distinction Lesson 6.3 established and Unit 7 resolves. Any group
that writes ``no solution'' for row 3 or row 4 should be corrected on the spot.
\end{teachernote}

\end{document}
